\section{Множество меры ноль по Лебегу. Свойства таких множеств. Критерий Лебега интегрируемости по Риману (без док).}

\subsection{Множество меры ноль по Лебегу.}
\begin{definition}
    Множество $E \subset \R$ называется множеством нулевой меры Лебега, если

    $\forall \eps > 0~\exists\Subst{\{(a_i; b_i)\}}{i = 1}{+\infty}$:\\
    1) $\Sum{i = 1}{+\infty}(b_i - a_i) \leq \eps$\\
    2) $E \subset \bigcup_{i=1}^{+\infty} (a_i, b_i)$

    Обз: $\lambda(E) = 0$\\
\end{definition}

\begin{example}
    $\forall$ счетное множество $\subset \R$ - меры ноль по Лебегу
    \begin{proof}
        $\Subst{\tau_n}{n = 1}{+\infty}$ - счетное множество

        Тогда нам нужно покрыть все точки из этого множества. Это можно сделать, взяв окрестности этих точек.

        $\Sum{i = 1}{+\infty}\frac{1}{2^{i+1}} = \frac{1}{2}$\\
        Такая сумма была взята, чтобы потом при сложении длин интервалов у нас получилось значение меньше эпсилон\\
        Стоит помнить, что окрестность точки в два раза меньше, чем ее радиус(ещё одна причина, по которой мы взяли такую сумму)\\ 
        $(a_i; b_i) = U_{\frac{\eps}{2^{i+1}}}(\tau_i)$\\
    \end{proof}
\end{example}

\subsection{Свойства таких множеств.}
\begin{definition}[Свойства]
    1) $\lambda(A) = 0$ и $B \subset A \implies \lambda(B) = 0$\\
    Док-во: возьмем все интервалы для А, они покроют В, так как В вложено в А\\
    2) $\lambda(A) = 0$ и $\lambda(B) = 0 \implies \lambda(A \cup B) = 0$\\
    Док-во: Возьмем $\eps$ и построим покрытие с суммой меньше $\frac{\eps}{2}$ для А и для B\\
    Если мы наложим эти покрытия, то сумма всех интервалов будет меньше $\eps$ и они точно покроют все объединение А и В. 
\end{definition}

\subsection{Критерий Лебега интегрируемости по Риману(без док)}

\begin{theorem}[Критерий Лебега интегрируемости по Риману]
    $f(x) \in R[a; b] \iff \text{множество точек разрыва - меры ноль по Лебегу}$
\end{theorem}